\chapter[Structure de l'espace des AC]%
{Structure de l'espace des automates cellulaires}
\mysetref{rule-space}

Intuitivement, la classification des~AC est une classification
qualitative de leurs comportements dynamiques. %
Cette classification repose sur des mesures statistiques des
comportements\footnote{%
  Par exemple l'entropie %
  $S = - \sum_{i} p_{i} \log(p_{i})$ %
  ou l'information mutuelle~(la corr\'elation de deux distributions) %
  $M = \sum_{i} \sum_{j} p_{ij} \log \frac{p_{ij}}{p_{i} p_{j}}$~\cite[]{Li90b}.}. %

La premi\`ere question est celle de l'existence d'un certain
comportement, une r\'eponse g\'en\'erale est donn\'ee par l'existence d'AC
universels. %
Un cadre g\'en\'eral est donn\'e par l'utilisation des m\'ethodes de la
m\'ecanique statistique pour qualifier les comportements. %

La deuxi\`eme question est celle du nombre d'AC d'un type donn\'e et de
leurs localisations dans l'espace des r\`egles. %
Cette question introduit le probl\`eme de l'existence de param\`etres
permettant de parcourir l'espace des r\`egles. %

Le probl\`eme de la classification des~AC se situe \`a l'intersection du
probl\`eme \aquote{forward}\footnote{%
  D\'eterminer les propri\'et\'es \`a partir de la r\`egle.} %
et inverse\footnote{%
  D\'eterminer la r\`egle \`a partir des propri\'et\'es.}~%
\cite[]{Gutowitz90a}. %
Il implique l'existence d'un ensemble de mesures r\'ealisables sur les
r\`egles et sur leurs dynamiques. %

L'existence de param\`etres de contr\^ole de la dynamique des~AC permet
de d\'eterminer une structure pour l'espace des~AC. %

%%
%% MINIMAL VERSION
%%
\section{Une version minimale du probl\`eme}

Les~AC \'el\'ementaires, tel que d\'efinis par~\cite{Wolfram86a}
sont les~AC de dimension~$d = 1$, de rayon\footnote{%
  Le nombre de cellules voisines est $d r + 1$.} %
$r = 1$ et de nombre d'\'etats~$k = 2$. %
Une simplification ultime consid\`ere l'espace des~AC de param\`etres %
$d = 0$, $r = 0$, $k = 2$, %
Ce qui donne $k^{^{d r + 1}} = 4$ r\`egles~(Z\'ero, Flip, Same, One),
lesquelles donnent 8~dynamiques:

\noindent Z(0) = 0*, Z(1) = 10*, O(0) = 01*, O(1) = 1*\\
S(0) = 0*, S(1) = 1*, F(0) = (01)*, F(1) = (10)*

soit 4~points fixes sans transitoire~(en 2 vari\'et\'es), 2 points fixes
avec transitoire de taille 1, 2 cycles de taille 2 sans transitoire.

Les~AC \'el\'ementaires, dont le nombre est~256 peuvent \^etre
\'enum\'er\'es\footnote{%
  Des donn\'ees statistiques peuvent \^etre obtenues pour
  chaque r\`egle appliqu\'ee sur un espace de configuration.}, %
mais l'\'etude des dynamiques associ\'ees \`a une r\`egle d\'epend du
nombre~$N$ de cellules de l'espace de configuration~(qui tend
vers l'infini) et fait n\'ecessairement appel \`a des m\'ethodes statistiques. %
Pour $d = 1$, $k = 2$, $r = 2$, le nombre de r\`egle est $2^{32}$ et ne
permet d\'ej\`a plus une \'enum\'eration.

\mytwoplotfull[htb]{lambda}{flip}%
{}%
{Les 128 premi\`ere valeur de \protect\alambda et \protect\Flip, pour un calcul sur les
  16384 premiers entier}%
{Le param\`etre appel\'e \alambda par~\cite{Li90b} est d\'efini par %
  $\lambda = \frac{k^{dr+1}-\#0}{k^{dr+1}}$~%
  \mygetref[paren]{lambda}. %
  Pour les r\`egles binaires, \alambda est une version normalis\'ee de la
  fonction $\#_{1}(n) = n - \sum_{i=1}\lfloor n/2^{i}\rfloor$~%
  \cite[]{Wolfram83} donnant le nombre de bits \`a~1 dans la
  repr\'esentation binaire d'un nombre. %
  Nous d\'efinissons \Flip pour les r\`egles binaires comme le nombre de
  \aconcept{bit-flip} rencontr\'es en balayant les entr\'ees de la table de
  transition~\mygetref[paren]{flip-par}. %
  Les calculs sont fait sur les 16384 premiers entiers consid\'er\'es comme
  autant de r\`egles d'un hypoth\'etique voisinage \`a~$\log_{2}(14)$
  voisins.\\
  $\#0$ d\'esigne le nombre de z\'eros de la table de transition.}

\mytwoplotfull[htb]{lambda-deriv}{flip-deriv}%
{}%
{Les 256 premi\`eres valeurs de la d\'eriv\'ee de \protect\alambda et \protect\Flip}%
{La d\'eriv\'ee est calcul\'ee \`a partir des donn\'ees pr\'esent\'ees en
  figure~\mygetplot{lambda+flip}:\\ %
  $f^{\prime}_{i} = f_{i} - f_{i-1}$.}

\mytwoplotfull[htb]{lambda-prob}{flip-prob}%
{}%
{Les 512 premi\`eres valeurs des fonctions \protect\alambda et
  \protect\Flip, et de leur probabilit\'e d'apparition}%
{La distribution des valeurs d'une fonction utilis\'ee comme para\-m\`e\-tre
  d'\'echantillonage d'un espace de r\`egles influence la nature de
  l'\'echan\-til\-lonage. %
  Les exp\'eriences utilisant \alambda comme param\`etre de
  contr\^ole~(\cite{Li90b} par exemple), ne tiennent pas compte de la
  non lin\'earit\'e de la distribution des probabilit\'es d'apparition des
  valeurs de \alambda sur un espace de r\`egles.}

\mytwoplotfull[htb]{lambda-versus-prob}{flip-versus-prob}%
{}%
{Plot de \protect\alambda et \protect\Flip \protect\vs leur probabilit\'e d'apparition}%
{Nous n'avons pas actuellement d'explication compl\`ete de ces
  r\'esultats, il est n\'eanmoins clair que \alambda et \Flip poss\`edent des
  propri\'et\'es suffisamment diff\'erentes pour justifier la poursuite
  d'exp\'eriences utilisant \Flip} 

\mytwoplotfull[htb]{lambda-flip-hist}{lambda-versus-flip}%
{}%
{L'histogramme des distribution de \protect\alambda et \protect\Flip, et
  \protect\alambda \protect\vs \protect\Flip.}%
{Les valeurs de l'histogramme apparaissent ici non normalis\'ees, \ie en
  nombre de bits \`a~1 pour \alambda et en nombre de bit-flips pour
  \Flip.}

\mytwoplotfull[htb]{lambda-flip-ratio-sorted}{lambda-flip-prob-sorted}%
{}%
{Un tri des valeurs de la fonction et leurs probabilit\'es d'apparition
  pour \protect\alambda et \protect\Flip}%
{Le nombre de paliers correspond au nombre de valeurs possibles pour
  chaque fonction, faible pour notre espace \`a voisinage
  de~$\log_{2}(14)$~bits.}

\myplotfullR[htb]{.75}{E-lambda-1}%
{}%
{Valeurs de \protect\alambda et \protect\En}%
{Nous avons calcul\'e les valeurs de \alambda et \En{}\footnote{%
    \aconcept{L'ench\^assement}, voir \'equation~\ref{rs-eq1}.} %
  pour un sous-espace \'echantillon de l'espace des r\`egles binaires \`a
  voisinage de Moore. %
  Le sous espace est constitu\'e des r\`egles additives, centre inclu, dont la table
  secondaire supporte les fonctions constantes \aquote{z\'ero} et
  \aquote{un} et la fonction \aquote{unchanged}. %
  L'espace contient donc $3^{10} = 59049$ r\`egles. %
  Les valeurs plot\'ees pour \En sont:~$\mEn_{centre}$,
  $\mEn_{rayon}$ et la valeur moyenne de \En sur tous les
  voisins. %
  Le pic montre la position de \arule{Anneal} dans le sous-espace de
  r\`egles. %
  Nous avons tent\'e de corr\'eler les param\`etres calcul\'es sur l'espace des
  r\`egles aux propri\'et\'es des trajectoires associ\'ees \`a ces r\`egles pour des
  configurations initiales al\'eatoires de densit\'e\,\Sundemi sur des espaces
  de tailles entre \asquare{64} et \asquare{256}, en vue de
  d\'eterminer les restrictions suppl\'ementaires \`a appliquer \`a une
  recherche automatique.}

%%
%% WOLFRAM CLASS
%%
\section{La classification de Wolfram}

La classification de Wolfram est une classification de r\'ef\'erence pour
de nombreux
auteurs~\cite[]{Vichniac90,Li90b,Wootters90,McIntosh90,Chate90,Gutowitz90b}.

Dans les quatre classes de Wolfram, %
les classes I, II et III sont vues comme des attracteurs de syst\`emes
dynamiques:

\subsection*{classe I}

\begin{itemize}
\item point limite
\item ou tendance \`a un \'etat spatialement homog\`ene
\item ou les formes disparaissent avec le temps
\item ou de petits changements sur la configuration initiale sont
  sans effet sur l'\'etat final
\end{itemize}

\subsection*{classe II}

\begin{itemize}
\item cycle limite
\item ou montrant une s\'equence de structure simple, stable ou
  p\'eriodique
\item ou les formes \'evoluent vers une taille fixe finie
\item ou de petits changements sur la configuration initiale
  induisent des changements sur des r\'egions de tailles finies
\end{itemize}

\subsection*{classe III}

\begin{itemize}
\item attracteur chaotique
\item ou montrant un comportement chaotique ap\'eriodique
\item ou les formes croissent ind\'efiniment \`a une vitesse fixe
\item ou de petits changements sur la configuration initiale
  induisent des changements sur des r\'egions de tailles toujours
  croissantes
\end{itemize}

\subsection*{classe IV}

\begin{itemize}
\item comportement plus complexe que les pr\'ec\'edents, conjecture de
  calcul universel
\item ou montrant des structures localis\'ees complexes,
  \'eventuellement propageantes
\item ou les formes croissent et se contractent avec le temps
\item ou de petits changements sur la configuration initiale
  induisent des changements irr\'eguliers
\end{itemize}

La description des classes varie selon les auteurs. %
On peut introduire la notion de d\'ecalage spatial
homog\`ene~\cite[]{Li90b}, %
souligner la difficult\'e que pr\'esente certains comportements
interm\'ediaires~\cite[]{McIntosh90}, %
ou proposer des interpr\'etations structurelles faisant, par exemple,
des~AC de la classe~IV un \emph{m\'elange} d'AC de la classe~II
et~III~\cite[]{Kayama93}. %
Les 4 classes~(fixe, p\'eriodique, chaotique, complexe) apparaissent
actuellement comme un point de d\'epart in\'evitable de toute tentative de
classification. %
La question de la structure de l'espace, et donc des param\`etres de
contr\^ole devient alors la question premi\`ere.

%%
%% LAMBDA
%%
\section{Le param\`etre $\lambda$ de Langton}
\label{ref-lambda}

Le param\`etre~$\lambda$ est d\'efini comme le pourcentage d'entr\'ees de la
table de transition dont la valeur est diff\'erente de l'\'etat z\'ero:~%
$\lambda = \frac{k^{dr+1}-\#0}{k^{dr+1}}$. %

A une petite variation de~$\lambda$~\cite[]{Langton92b} entre deux
r\`egles correspond une petite distance de Hamming entre deux r\`egles. %
La variation du param\`etre~$\lambda$ permet d'\'echantillonner
un espace de r\`egles pour~$d$, $k$ et~$r$ fix\'es. %
Chaque r\`egle est appliqu\'ee sur un ensemble de configurations
initiales, et des mesures statistiques sont faites \`a la vol\'ee. %
Les s\'eries temporelles de mesures sont utilis\'ees pour d\'eterminer
l'appartenance de la r\`egle \`a une classe ou une autre. %
On peut alors essayer d'\'etablir une corr\'elation entre~$\lambda$
et la classe d'appartenance de l'AC. %

Cette proc\'edure est a priori applicable pour d'autres param\`etres, %
et vise \`a travers l'\'echantillonnage \`a r\'ev\'eler la structure g\'en\'erale
que la variation du param\`etre de contr\^ole imprime \`a l'espace des r\`egles. %

Une question associ\'ee \`a la structure de l'espace d\'efini par un
param\`etre est la pr\'esence et le type d'une \'eventuelle transition de
phase, qui associe \`a une valeur critique du param\`etre de contr\^ole une
transition plus ou moins brusque de la classe de comportement des
r\`egles. %

Une autre question est la pr\'esence d'une tendance au continu associ\'ee
\`a l'augmentation de la taille de l'espace des r\`egles. %

%%
%% EXPERIMENTS
%%
\section[param\'etrage des classes de comportements]%
{Exp\'erimentations sur le param\'etrage des classes de comportements}

Li, Packard et Langton~(\citeyear{Li90b}) proposent l'existence de
transitions de phase contr\^ol\'ees par le param\`etre~$\lambda$. %
Il s'agit d'une transition du premier ordre~(transition brutale) entre
les r\`egles de type II et III, %
ou d'une transition de deuxi\`eme ordre~(transition graduelle) associant
les r\`egles de type IV \`a une zone d'intersection des types II et III\@. %
Les mesures statistiques utilis\'ees sont l'entropie, l'information
mutuelle et la vitesse de propagation des diff\'erences sur l'espace
cellulaire\footnote{%
  Consid\'er\'ee comme une mesure de l'impr\'edictibilit\'e de l'\'evolution
  des \emph{formes} de l'espace cellulaire, elle devient une mesure
  de la qualit\'e chaotique du comportement de l'AC.}. %
Les espaces de r\`egles sont typiquement born\'es par $d <= 2$, $k < 10$,
$r < 12$. %
Ce type d'exp\'eriences a d'abord \'et\'e men\'e sur les~AC
\'el\'ementaires~\cite[]{Li90a}. %
\cite{Wootters90} \'etendent la recherche \`a des
espaces de r\`egles plus vastes, dans le but de r\'eduire la dispersion de
la valeur critique. %
L'augmentation de~$r$ supprime la transition~(la valeur critique tend
vers z\'ero), l'augmentation de~$k$ tend \`a r\'eduire la dispersion de la
valeur critique.

Toutefois, la critique d'une exp\'erience de \cite{Packard88},
formul\'ee par \cite{Mitchell93a} fait appara\^itre les limites
de l'hypoth\`ese de Langton sur l'\'emergence du \emph{calcul}\footnote{%
  Il s'agit en fait de la transition vers des r\`egles dont les mesures
  statistiques indiquent une appartenance \`a la classe IV, lesquelles
  sont suppos\'ees universelles.} %
\`a la fronti\`ere du chaos~\cite[]{Langton90,Langton92b}. %
Packard utilise un algorithme g\'en\'etique pour faire \'evoluer des r\`egles
effectuant un calcul particulier, et interpr\`ete les r\'esultats comme
une indication de l'existence d'une pression s\'elective sur l'\'evolution
des r\`egles vers les valeurs critiques de~$\lambda$ associ\'ees \`a une
transition de phase et \`a la r\'ealisation d'un calcul complexe.

%% 0.27 quoi
Mitchell trouve des r\'esultats diff\'erents, et souligne la n\'ecessit\'e
d'une meilleure compr\'ehension~(d\'efinition) des concepts utilis\'es. %
Mitchell montre, par exemple, que l'\'evolution par algorithme g\'en\'etique
d'une r\`egle calculant l'\'etat majoritaire de la configuration initiale
en tendant vers un point fixe homog\`ene de l'\'etat majoritaire, %
se trouve n\'ecessairement loin du point critique~(0.27 d'apr\`es Wooter
et Langton) avec $\lambda$ = 0.5. %

Li a par ailleurs sugg\'er\'e~\cite[]{Li89}, en apparente contradiction
avec l'id\'ee d'une association entre calcul, complexit\'e et ph\'enom\`enes
critiques~\cite[]{Li90b} une structure fractale~\cite[]{Mandelbrot84}
pour la distribution des r\`egles complexes dans l'espace des r\`egles.

%%
%% PROPOSITIONS
%%
\section[Le param\'etrage d'ench\^assement]%
{Proposition pour une recherche d'un param\'etrage de l'ench\^as\-sement}
\mysetref{flip-par}

Face aux probl\`emes soulev\'es par l'utilisation du param\`etre~$\lambda$,
nous proposons quelques remarques pouvant servir de point de d\'epart \`a
de nouvelles exp\'eriences sur la structure de l'espace des~AC.

La premi\`ere remarque concerne l'existence d'autres param\`etres
de contr\^oles que~$\lambda$. %
Li, Packard et Langton~\cite[]{Li90b} soulignent l'importance de la
recherche d'autres param\`etres de contr\^ole et sugg\`erent que la
dispersion de la valeur critique pourrait \^etre r\'eduite par un
param\`etre suppl\'ementaire~(le \emph{param\`etre myst\'erieux} permettant de
localiser des fronti\`eres plus pr\'ecises entre les diff\'erents comportements). %
Il nous semble n\'eanmoins peu probable qu'un seul param\`etre
suppl\'ementaire suffise.

La deuxi\`eme remarque concerne la structure ench\^ass\'ee des espaces de
r\`egles. %
Nous sugg\'erons que l'\'etude du degr\'e d'ench\^assement d'une r\`egle dans un
espace permet la d\'efinition d'un nouveau param\`etre de contr\^ole.

La troisi\`eme remarque concerne les techniques de r\'eductions du nombre
de r\`egles l\'egales \`a consid\'erer pour un espace donn\'e, c'est \`a dire
l'ablation des r\`egles anisotropiques ou ench\^ass\'ees, et la possibilit\'e
d'utiliser un mapping interm\'ediaire.

\subsection*{Le param\`etre Flip}

Nous proposons l'utilisation du param\`etre~\emph{Flip} d\'efini comme le
nombre de changements d'\'etats rencontr\'e en balayant successivement
chaque entr\'ee de la table de transition.

\[{\rm flip} =
\frac{1}{N}\sum_{i=0}^{N-1}{\rm diff}(f_{i}-f_{(i+1){\rm mod}N})\]

avec $N=k^{dr+1}$, et

\[{\rm diff}(x) = \left\lbrace
\begin{array}{l}
  1\ {\rm if}\ x \ne 0\\ 0\ {\rm if}\ x = 0
\end{array}
\right.\]

Si l'on plonge\footnote{%
  Nous d\'efinissons le plongement de r\`egles comme augmentation de la
  taille du voisinage sans influence sur les trajectoires.} %
une r\`egle~$R_{N}$ param\'etr\'ee par~$r = N$ dans une
r\`egle~$R_{N + 1}$, l'insensibilit\'e de la r\`egle aux deux voisins
suppl\'ementaires se traduit par une duplication des entr\'ees de la r\`egle
$R_{N}$. %
Par exemple pour~$k = 2$, si pour la colonne~$i$ l'\'egalit\'e de chacune
des~$N/2$ paires~$x_{j} = x_{j + 2^{i}}$ est v\'erifi\'ee, la colonne~$i$
est inutilis\'ee, et soit la r\`egle est anisotropique, soit la r\`egle est
isomorphe \`a une r\`egle d'espace de taille inf\'erieure et peut \^etre
isotropique.

La somme des degr\'es d'\'egalit\'e de chaque colonne est un candidat
possible au r\^ole de param\`etre de contr\^ole.

Parce que l'on peut construire une r\`egle pour un espace de taille
donn\'ee en plongeant une r\`egle d'un espace de taille inf\'erieure~(par
duplication des entr\'ees), \`a~$\lambda$ et Flip constant, puis d\'eplacer
peu \`a peu la r\`egle vers l'espace de taille sup\'erieure en flippant des
entr\'ees~(\'eventuellement en conservant~$\lambda$ constant), la mesure
de Flip pourrait \^etre associ\'ee au degr\'e d'ench\^assement d'une r\`egle.

L'utilisation d'un mapping interm\'ediaire permet de r\'eduire la taille
de l'espace des r\`egles, tout en maintenant une coh\'erence forte sur le
sous-ensemble choisi. %
Par exemple le r\'eduction du voisinage \`a une somme excluant la cellule
centrale, conserve un sous ensemble isotropique de r\`egle. %
Pour~$d = 2$, $k = 2$, $r = 1$ les~$2^{512}$ r\`egles
deviennent~$4^{9}$, en utilisant les 4 op\'erations des r\`egles
minimales~(Z\'ero, Same, Flip, One) comme valeur de la table
interm\'ediaire. %

\subsection*{Le param\`etre d'ench\^assement}

Nous d\'efinissons le param\`etre d'ench\^assement d'un voisin pour une
r\`egle comme l'influence de ce voisin sur la r\`egle. %
Soit une r\`egle $R$ d'un AC de $k = 2^n$ \'etats, $d$ dimensions et de
rayon de voisinage $r$. %
$N = k^{dr+1} = 2^b$ est le nombre d'entr\'ees de la table. %
Pour un AC de $k = 2$ \'etats les $b = dr+1$ bits du voisinage donnent
un vecteur $E$ de $b$ valeurs.

%%\begin{eqnarray}
%%\label{rs-eq1}
%%        E_{b} = \frac{1}{N} \sum_{i = 0}^{N/2 - 1} R_{e_{ib}} - R_{e_{ib} + 2^b} \\
%%        \mbox{avec : } e_{ib} = \frac{i}{2^b} 2^{b + 1} + i \bmod 2^b
%%        = 2i + i \bmod 2^b \nonumber
%%\end{eqnarray}
\begin{eqnarray}
\label{rs-eq1}
        E_{b} = \frac{1}{N} \sum_{i = 0}^{N/2 - 1} R_{e_{ib}} - R_{e_{ib} + 2^b} \\
        \mbox{avec : } e_{ib} = 2i + i \bmod 2^b \nonumber
\end{eqnarray}

Une valeur moyenne $E = \frac{1}{b} \sum_{i = 0}^{b} E_b$ peut \^etre
d\'efinie, pour servir de param\`etre de contr\^ole \`a la fa\c{c}on du
param\`etre $\lambda$ cependant d'autre r\'eduction que la moyenne
pourrait \^etre choisie, afin d'exploiter l'information sur le degr\'e de
sym\'etrie $R$.

\begin{eqnarray}
        E_{life} & = &
        {0.617\ 0.617\ 0.617\ 0.617\ 0.890\ 0.617\ 0.617\ 0.617\ 0.617} \nonumber \\
        \hat{E}_{life} & = & 0.647569 \nonumber \\
        \lambda_{life} & = & 0.273438 \nonumber
\end{eqnarray}

Le probl\`eme de la structure de l'espace des~AC, semble \^etre un
probl\`eme encore largement ouvert. %
La question des formes possibles de la dynamique des~AC pour des
espaces de r\`egles de grande taille, et la question d'une combinatoire
de r\`egles permettant de contr\^oler l'\'emergence de ces formes ne
poss\`edent pas actuellement de r\'eponses.

%%
%% Local Variables:
%% mode: auto-fill
%% iso-accents-minor-mode: nil
%% TeX-master: "top"
%% TeX-command-default: "mytex"
%% Local IspellParsing: "latex-mode ~latin1"
%% LocalWords: "forward Wolfram Flip Same One transient CLASS II III IV LAMBDA"
%% LocalWords: "Langton Hamming l'AC EXPERIMENTS Li Packard Wootters Mitchell AC"
%% LocalWords: "Wooter mapping flip flippant em d'AC spatialement ap\'eriodique"
%% LocalWords: "propageantes param\'etrage l'impr\'edictibilit\'e fractale l'ench\^as"
%% LocalWords: "sement anisotropiques anisotropique isotropique rule-space htb"
%% LocalWords: "lambda paren flip-par lambda-deriv flip-deriv lambda-prob versus"
%% LocalWords: "flip-prob lambda-versus-prob flip-versus-prob lambda-flip-hist"
%% LocalWords: "lambda-versus-flip lambda-flip-ratio-sorted E-lambda sous-espace"
%% LocalWords: "lambda-flip-prob-sorted Moore unchanged plot\'e touts Anneal"
%% End:
